\chapter{Elementos Adicionais de Texto e Estrutura Semântica em HTML}

\section{Citações em HTML}

Quando você utiliza uma ideia que não é sua, é fundamental citar a fonte — tanto por uma questão ética quanto, em muitos contextos, por uma questão legal. O HTML oferece elementos específicos para marcar citações.

\subsection{Citação em Bloco: o Elemento blockquote}

O elemento blockquote é usado quando você cita, na íntegra, um conteúdo mais extenso — um trecho completo de outra fonte, e não apenas uma palavra ou frase isolada.

\begin{codebox}{HTML}
 <p>Segundo a documentação da Mozilla:</p>
 <blockquote cite="https://developer.mozilla.org/pt-BR/">
   <p>A MDN Web Docs é um recurso aberto e colaborativo para
      desenvolvedores, por desenvolvedores, documentando tecnologias
      web, incluindo CSS, HTML e JavaScript.</p>
 </blockquote>
\end{codebox}

O atributo cite indica de onde a citação foi extraída. É importante notar que esse atributo não é visível na tela — ele existe apenas como metadado, útil caso você (ou algum script) queira acessá-lo programaticamente.

\subsection{Citação em Linha: o Elemento q}

Quando a citação é apenas um trecho dentro de um parágrafo maior — e não o parágrafo inteiro — utiliza-se o elemento q (quotation), que é um elemento em linha (inline).

\begin{codebox}{HTML}
<p>
  De acordo com a documentação da MDN, um dos objetivos do projeto é
  <q cite="https://developer.mozilla.org/">fornecer um recurso completo
  e confiável para desenvolvedores web de todos os níveis</q>.
</p>
\end{codebox}

Muitos navegadores adicionam automaticamente aspas ao redor do conteúdo de um elemento q, mesmo sem que você as digite manualmente — outra vantagem de usar marcação semântica ao invés de digitar as aspas você mesmo.

\subsection{O Elemento cite}

Existe ainda um terceiro elemento relacionado a citações, o cite, que não deve ser confundido com o atributo cite de blockquote e q. O elemento cite é usado para indicar o título de uma obra — um livro, um artigo, uma música.

\begin{infobox}{Dica}
Nenhum desses três recursos (elemento blockquote, elemento q ou atributo cite) torna a referência clicável automaticamente. Uma boa prática recomendada é envolver o elemento cite (ou o texto da fonte) dentro de um elemento de âncora (a), permitindo que o leitor acesse a fonte original com um clique:
\end{infobox}

\begin{codebox}{HTML}
<blockquote cite="https://developer.mozilla.org/">
  <p>Este é um exemplo de citação em bloco.</p>
</blockquote>
<p>
  — <a href="https://developer.mozilla.org/">
      <cite>MDN Web Docs</cite>
    </a>
</p>
\end{codebox}

\section{Abreviações, Endereços e Outras Marcações Textuais}

\subsection{Abreviações: o Elemento abbr}

O elemento abbr marca abreviações e siglas, utilizando o atributo title para fornecer a forma expandida do termo. A vantagem prática é que, ao passar o mouse sobre a sigla, o navegador exibe automaticamente um tooltip com o significado completo.

\begin{codebox}{HTML}
<p>
  Nós usamos <abbr title="HyperText Markup Language">HTML</abbr>
  para estruturar nossos documentos web.



 </p>
\end{codebox}

\begin{infobox}{Dica}
O HTML já teve um elemento acronym, específico para acrônimos (siglas que se pronunciam como palavras). Esse elemento se tornou obsoleto e não é mais suportado de forma consistente pelos navegadores modernos — utilize sempre abbr, tanto para siglas quanto para acrônimos.
\end{infobox}

\subsection{Endereços: o Elemento address}

O elemento address marca um trecho de texto que representa informações de contato — endereço residencial, comercial ou de correspondência de uma pessoa ou organização.

\begin{codebox}{HTML}
 <address>
   Maria Oliveira<br>
   Rua das Palmeiras, 456<br>
   Manchester, Reino Unido
 </address>
\end{codebox}

Repare no uso do elemento br (line break), que insere uma quebra de linha simples — lembre-se de que o HTML ignora espaços em branco e quebras de linha do códigofonte, então, sem o br, todo esse texto apareceria em uma única linha corrida.

\begin{infobox}{Dica}
O elemento address deve conter apenas informações de contato relacionadas ao autor do documento (ou da seção em que está inserido) — não o utilize para qualquer endereço genérico mencionado no texto, como o endereço de uma empresa citada em uma notícia.
\end{infobox}

\subsection{Superescrito e Subescrito}

Os elementos sup (superescrito) e sub (subescrito) posicionam o conteúdo ligeiramente acima ou abaixo da linha do texto, respectivamente — úteis para datas por extenso, fórmulas químicas e expressões matemáticas.

\begin{codebox}{HTML}
 <p>Meu aniversário é no dia 25<sup>th</sup> de maio.</p>
 <p>A fórmula química da cafeína é 
   C<sub>8</sub>H<sub>10</sub>N<sub>4</sub>O<sub>2</sub>.<
 <p>Considere a equação x<sup>2</sup> + 2x + 1 = 0.</p>
\end{codebox}

\subsection{Representação de Código-Fonte}

Ao escrever documentação técnica, é comum precisar representar código-fonte, comandos de teclado e saídas de terminal. O HTML oferece elementos específicos para cada um desses casos: code (código genérico), var (variáveis), kbd (keyboard, entradas de teclado) e samp (sample output, saída de um programa).

\begin{codebox}{HTML}
<pre><code>
let x = 42;
console.log(x);
</code></pre>

<p>
  A variável <var>x</var> armazena o valor calculado.
</p>

<p>
  Para copiar, pressione <kbd>Ctrl</kbd> + <kbd>C</kbd>.
</p>

<p>
  A saída do comando foi: <samp>Build concluído com sucesso.</samp>
</p>
\end{codebox}

Note o uso do elemento pre (preformatted text) envolvendo o code: diferentemente do restante do HTML, o pre preserva espaços em branco e quebras de linha exatamente como foram digitados no código-fonte — essencial para exibir trechos de código com indentação correta.

\subsection{Datas e Horários: o Elemento time}

Como sites e aplicações web são utilizados no mundo inteiro, com formatos de data distintos entre culturas, o elemento time permite indicar, de forma padronizada e não ambígua, o valor de uma data ou horário através do atributo datetime, enquanto o conteúdo visível do elemento pode ser escrito da forma mais amigável para o leitor humano.

\begin{codebox}{HTML}
<p>
  Publicado em <time datetime="2016-01-20">20 de janeiro de 2016</time>.
</p>
<p>
  O evento começa às <time datetime="2016-01-20T19:30">19h30</time>.
</p>
\end{codebox}

\section{Estrutura Semântica de Documentos}

Até aqui, vimos elementos que marcam trechos específicos de texto dentro de um parágrafo. Agora avançamos para um nível mais amplo: os elementos que estruturam as grandes áreas de uma página web — cabeçalho, navegação, conteúdo principal, barra lateral e rodapé. Todo site, independentemente do seu propósito, costuma compartilhar essas mesmas grandes divisões. A figura a seguir ilustra visualmente como essas tags se organizam em uma página típica.

\begin{infobox}{header}
logotipo, título do site
\end{infobox}

\begin{infobox}{nav}
menu de navegação principal
\end{infobox}

\begin{infobox}{main}
article / section conteúdo principal aside barra lateral, article / section conteúdo relacionado mais conteúdo
\end{infobox}

\begin{infobox}{footer}
direitos autorais, contato
\end{infobox}

\subsection{O Elemento header}

O elemento header representa um cabeçalho introdutório, que normalmente contém o logotipo, o título do site ou da seção, e eventualmente um menu de navegação ou um campo de busca. É importante entender que header pode aparecer mais de uma vez no documento: um header diretamente dentro do body representa o cabeçalho do site inteiro, mas um header dentro de um article ou section representa o cabeçalho específico daquela seção.

\begin{codebox}{HTML}
<header>
  <h1>Minha Empresa</h1>
  <img src="logo.png" alt="Logotipo da Minha Empresa">
  <nav>
    <ul>
      <li><a href="/">Início</a></li>
      <li><a href="/sobre">Sobre</a></li>
      <li><a href="/contato">Contato</a></li>
    </ul>
  </nav>
</header>
\end{codebox}

\subsection{O Elemento nav}

O elemento nav agrupa os links de navegação principal do site. É importante reserválo apenas para os menus mais importantes — links secundários, presentes apenas em rodapés ou em contextos específicos, não precisam necessariamente estar dentro de um nav.

\subsection{O Elemento main}

O elemento main envolve o conteúdo principal e único daquela página específica — ou seja, o conteúdo que não se repete em outras páginas do site (diferentemente do cabeçalho e do rodapé, que normalmente são os mesmos em todas as páginas). Deve haver, no máximo, um elemento main por página.

\subsection{Os Elementos article e section}

Dentro do main, o conteúdo costuma ser subdividido em article e/ou section. A distinção entre eles não é rígida e depende do contexto:

\begin{itemize}
  \item article representa um conteúdo autocontido e independente, que faria sentido mesmo se distribuído isoladamente (por exemplo, uma postagem de blog, uma notícia, um comentário de usuário).
\end{itemize}

\begin{itemize}
  \item section agrupa um conjunto de conteúdo tematicamente relacionado, mas que não necessariamente é independente — por exemplo, uma seção sobre ”música”e, mais abaixo na mesma página, outra seção sobre ”culinária”.
\end{itemize}

\begin{codebox}{HTML}
<main>
  <article>
    <header>
       <h2>Como aprender CSS do zero</h2>
    </header>
    <p>Neste artigo, vamos abordar os fundamentos do CSS...</p>
  </article>

  <section>
    <h2>Comentários dos leitores</h2>
    <article>
      <p>Ótimo artigo, me ajudou muito!</p>
    </article>
  </section>
</main>
\end{codebox}

\begin{infobox}{Dica}
Uma boa prática é sempre iniciar cada section (ou article) com um elemento de título (h1 a h6), da mesma forma que fizemos no capítulo anterior ao estruturar capítulos e subseções de um texto.
\end{infobox}

\subsection{O Elemento aside}

O elemento aside representa um conteúdo que está relacionado ao conteúdo principal, mas que não é o foco central da página — uma barra lateral, uma nota explicativa, um bloco de anúncios ou links relacionados.

\begin{codebox}{HTML}
<aside>
  <h2>Artigos relacionados</h2>
  <ul>
    <li><a href="/artigo-1">Introdução ao HTML</a></li>
    <li><a href="/artigo-2">Primeiros passos com CSS</a></li>
  </ul>
</aside>
\end{codebox}

\subsection{O Elemento footer}

O elemento footer representa o rodapé — normalmente contendo informações de contato, direitos autorais, links de política de acessibilidade ou de privacidade. Assim como o header, ele pode aparecer tanto em nível global (rodapé do site inteiro) quanto dentro de uma section ou article específico.

\begin{codebox}{HTML}
<footer>
  <p>&copy; 2026 Minha Empresa. Todos os direitos reservados.</p>
  <p><a href="/politica-privacidade">Política de Privacidade</a></p>
</footer>
\end{codebox}

\section{Um Exemplo Completo e a Questão da Acessibilidade}

Juntando todos esses elementos, uma página típica pode ser esboçada da seguinte forma:

\begin{codebox}{HTML}
<body>
  <header>
    <h1>Blog de Tecnologia</h1>
    <nav>
       <ul>
         <li><a href="/">Início</a></li>
         <li><a href="/artigos">Artigos</a></li>
       </ul>
    </nav>
  </header>

   <main>
     <article>
       <h2>Entendendo o Box Model do CSS</h2>
       <p>O modelo de caixa é um dos conceitos mais importantes...</p>
     </article>

     <aside>



         <h2>Sobre o autor</h2>
         <p>Desenvolvedor front-end apaixonado por acessibilidade.</p>
       </aside>
     </main>

  <footer>
    <p>&copy; 2026 - Todos os direitos reservados.</p>
  </footer>
</body>
\end{codebox}

Vale destacar por que essa estruturação semântica importa tanto: estima-se que entre 4\% e 5\% da população mundial tenha algum tipo de deficiência visual — e, considerando outras formas de deficiência, esse número pode chegar a cerca de 15\% da população global. Ao usar elementos genéricos como div e span em vez dos elementos semânticos apropriados, o desenvolvedor perde a oportunidade de comunicar corretamente o significado de cada parte do documento para leitores de tela e outras tecnologias assistivas.

\subsection{Quando Usar div e span}

Nem sempre existe um elemento HTML semântico adequado para o que se deseja representar — afinal, o HTML precisa ser genérico o suficiente para servir a qualquer tipo de site, e não conhece conceitos de nicho como ”carrinho de compras”. Nesses casos, é aceitável usar os elementos genéricos e não semânticos div (para blocos) e span (para conteúdo em linha), em conjunto com o atributo class para que o CSS ou o JavaScript possam localizá-los.

\begin{codebox}{HTML}
<div class="carrinho-compras">
  <ul>
    <li>Produto A - R$ 29,90</li>
    <li>Produto B - R$ 59,90</li>
  </ul>
  <p>Total: R$ 89,80</p>
</div>
\end{codebox}

\begin{infobox}{Dica}
O problema não é usar div e span em si — é usá-los sem perceber que eles não carregam nenhum poder semântico, quando existiria um elemento mais apropriado disponível. Sempre pergunte primeiro: ”existe um elemento HTML semântico que descreve melhor este conteúdo?”Só recorra a div/span quando a resposta for não.
\end{infobox}

Síntese do Capítulo

\begin{itemize}
  \item Citações usam blockquote (em bloco) ou q (em linha), com o atributo cite indicando a fonte; o elemento cite é reservado para títulos de obras.
\end{itemize}

\begin{itemize}
  \item Elementos como abbr, address, sup/sub, code/var/kbd/samp e time cobrem casos textuais específicos com forte valor semântico e de acessibilidade.
\end{itemize}

\begin{itemize}
  \item A estrutura de uma página web é organizada nos elementos header, nav, main, article, section, aside e footer.
\end{itemize}

\begin{itemize}
  \item header e footer podem existir tanto em nível global (o cabeçalho/rodapé do site) quanto em nível local (dentro de uma section ou article específico).
\end{itemize}

\begin{itemize}
  \item A diferença entre article e section não é rígida: article é conteúdo autocontido e independente; section agrupa conteúdo tematicamente relacionado.
\end{itemize}

\begin{itemize}
  \item div e span são elementos não semânticos, úteis apenas quando não existe uma alternativa semântica adequada disponível em HTML.
\end{itemize}

\begin{itemize}
  \item A escolha correta de elementos semânticos impacta diretamente a acessibilidade do site para leitores de tela e outras tecnologias assistivas.
\end{itemize}
